\documentclass[11pt]{article}
\usepackage{amsmath,amssymb,bm}
\usepackage[margin=2.5cm]{geometry}
\usepackage{booktabs}

\title{Two ways to make a model's output obey the rules\\
\large RAYEN and HardNet}
\author{}
\date{}

\begin{document}
\maketitle

\section{The problem}

At one instant the model produces a vector
\[
\mathbf{v} = (v_1,\;v_2,\;v_3,\;v_4,\;v_5,\;v_6).
\]

\paragraph{Rule 1 --- supply equals demand.} Write the signs as
\[
\mathbf{s} = (1,\;1,\;1,\;1,\;-1,\;1),
\]
and let $d$ be the known net demand. Then
\[
\boxed{\;\mathbf{s}\cdot\mathbf{P} \;=\; s_1P_1+s_2P_2+\dots+s_6P_6 \;=\; d\;}
\]

\paragraph{Rule 2 --- ramp rate}
\[
\boxed{\;\ell_i \;\le\; P_i \;\le\; u_i \quad\text{for every } i\;}
\]

Write the two rules as sets:
\[
C_0 = \{\mathbf{x} : \mathbf{s}\cdot\mathbf{x} = d\}
\quad\text{(a plane)},
\qquad
B = \{\mathbf{x} : \bm{\ell}\le\mathbf{x}\le\mathbf{u}\}
\quad\text{(a box)},
\qquad
C = C_0\cap B .
\]

Note $s_i^2=1$ for every $i$, so
\[
\|\mathbf{s}\|^2 = \mathbf{s}\cdot\mathbf{s} = s_1^2+\dots+s_6^2 = 6 = n. \tag{$\ast$}
\]

\section{What ``projection'' means}

The \textbf{projection} of a point $\mathbf{v}$ onto a set $C$ is the closest
point of $C$ to $\mathbf{v}$:
\[
\Pi_C(\mathbf{v}) \;=\; \arg\min_{\mathbf{P}\in C} \|\mathbf{P}-\mathbf{v}\|^2 .
\]

This is the one idea both methods are built from, but they apply it to different
objects:

\begin{center}
\begin{tabular}{ll}
\toprule
\textbf{HardNet} & projects the \emph{point} $\mathbf{v}$ onto $C$. Its output \emph{is} $\Pi_C(\mathbf{v})$.\\
\textbf{RAYEN} & projects a \emph{direction} onto the plane's parallel directions,\\
 & then travels along it from an anchor. Its output is \emph{not} $\Pi_C(\mathbf{v})$.\\
\bottomrule
\end{tabular}
\end{center}

\subsection*{Projection onto the plane $C_0$}

\begin{quote}
\textbf{Claim.} $\displaystyle \Pi_{C_0}(\mathbf{v}) = \mathbf{v} + \lambda\,\mathbf{s}$
with $\displaystyle \lambda = \frac{d-\mathbf{s}\cdot\mathbf{v}}{\|\mathbf{s}\|^2}
= \frac{d-\mathbf{s}\cdot\mathbf{v}}{6}$.
\end{quote}

\noindent\textbf{Proof.} Write the answer as $\mathbf{P}=\mathbf{v}+\mathbf{w}$, so
$\mathbf{w}$ is the correction we are choosing. Rule 1 says
\[
\mathbf{s}\cdot(\mathbf{v}+\mathbf{w}) = d
\quad\Longleftrightarrow\quad
\mathbf{s}\cdot\mathbf{w} = c, \qquad c := d-\mathbf{s}\cdot\mathbf{v}.
\]
So the task is: \emph{among all $\mathbf{w}$ with $\mathbf{s}\cdot\mathbf{w}=c$,
find the shortest.}

Split $\mathbf{w}$ into a part along $\mathbf{s}$ and a part perpendicular to it:
\[
\mathbf{w} = \lambda\,\mathbf{s} + \mathbf{z},
\qquad
\lambda := \frac{\mathbf{s}\cdot\mathbf{w}}{\|\mathbf{s}\|^{2}},
\qquad
\mathbf{z} := \mathbf{w}-\lambda\mathbf{s}.
\]
This split always works, and $\mathbf{z}$ really is perpendicular:
\[
\mathbf{s}\cdot\mathbf{z}
= \mathbf{s}\cdot\mathbf{w} - \lambda\,\|\mathbf{s}\|^{2}
= \mathbf{s}\cdot\mathbf{w} - \mathbf{s}\cdot\mathbf{w} = 0 .
\]

Now use the split twice.

\emph{First, the constraint fixes $\lambda$.} Taking the dot product with $\mathbf{s}$,
\[
c = \mathbf{s}\cdot\mathbf{w}
= \lambda\,\|\mathbf{s}\|^{2} + \underbrace{\mathbf{s}\cdot\mathbf{z}}_{=\,0}
= \lambda\,\|\mathbf{s}\|^{2}
\quad\Longrightarrow\quad
\lambda = \frac{c}{\|\mathbf{s}\|^{2}} .
\]
There is no choice about $\lambda$ --- every admissible $\mathbf{w}$ has this same one.

\emph{Second, Pythagoras kills $\mathbf{z}$.} Because $\mathbf{s}\perp\mathbf{z}$,
the cross term vanishes:
\[
\|\mathbf{w}\|^{2}
= \|\lambda\mathbf{s}+\mathbf{z}\|^{2}
= \lambda^{2}\|\mathbf{s}\|^{2} + 2\lambda\underbrace{(\mathbf{s}\cdot\mathbf{z})}_{=\,0} + \|\mathbf{z}\|^{2}
= \underbrace{\lambda^{2}\|\mathbf{s}\|^{2}}_{\text{fixed}} + \|\mathbf{z}\|^{2}.
\]
The first term cannot be changed. So $\|\mathbf{w}\|^2$ is smallest exactly when
$\mathbf{z}=\mathbf{0}$, giving $\mathbf{w}=\lambda\mathbf{s}$ and
\[
\boxed{\;\Pi_{C_0}(\mathbf{v}) = \mathbf{v} + \frac{d-\mathbf{s}\cdot\mathbf{v}}{6}\,\mathbf{s}\;}
\tag*{$\blacksquare$}
\]

\paragraph{So that is why $\mathbf{P}=\mathbf{v}+\lambda\mathbf{s}$.} Not because
``straight at the plane is shortest'' as an assertion, but because any correction
splits into a part along $\mathbf{s}$ that the constraint \emph{forces} on us and
a perpendicular part that only adds length. Throw the perpendicular part away.

\paragraph{In plain terms.} Coordinate $i$ changes by $\lambda s_i$, and
$|\lambda s_i| = |\lambda|$, so \emph{every generator moves by the same number of
megawatts}, up or down according to its sign.

\section{Approach A --- RAYEN}

\subsection*{Step 1: the anchor}

RAYEN travels \emph{from} somewhere, so it needs a starting point $\mathbf{a}\in C$
that already obeys both rules. HardNet needs no such thing. Where does
$\mathbf{a}$ come from?

\paragraph{The setting.} We are filling a gap of $N$ steps. We know the last real
value before the gap, $\mathbf{p}_L$, and the first real value after it,
$\mathbf{p}_R$. We fill the gap one step at a time; at step $j$ the previous
answer $\mathbf{P}_{\text{prev}}$ is already fixed, and
\[
k = N - j
\]
is the number of gap steps still to come. Counting the final move onto
$\mathbf{p}_R$, there are $k+1$ moves left.

The anchor is built in two stages: \textbf{glide}, then \textbf{snap}.

\subsubsection*{Stage 1 --- the glide (this is the interpolation)}

Draw a straight line from where we are to where we must end up, cut it into the
$k+1$ moves we have left, and take one:
\[
\boxed{\;\mathbf{c} = \operatorname{clip}\!\left(\mathbf{P}_{\text{prev}}
 + \frac{\mathbf{p}_R - \mathbf{P}_{\text{prev}}}{k+1},\ \bm{\ell},\ \mathbf{u}\right)}
\]
That fraction is the whole idea. If you are $60$ MW away with $3$ moves left, you
go $20$ MW. Next step you are $40$ away with $2$ left, so $20$ again --- a straight
line, arrived at by dividing the remaining distance by the remaining moves.

\paragraph{Why it is recomputed every step.} It would be equivalent to draw one
line at the start and follow it, \emph{if nothing ever pushed us off}. But Steps
2--6 below deliberately push us off: the ray moves the answer away from
$\mathbf{c}$. Redrawing the line from wherever we actually landed means the glide
always re-aims at $\mathbf{p}_R$ instead of running parallel to a line we have
already left. It cannot accumulate drift.

\paragraph{Why it never runs out of room.} The interval $[\bm{\ell},\mathbf{u}]$
carries a third term besides the ramp and the cap:
$\ell_i \ge p_{R,i} - k\,\text{rup}_i$ and $u_i \le p_{R,i} + k\,\text{rdn}_i$.
Read it as: \emph{never walk somewhere you can no longer get back from in the
$k$ steps you have left}. With that guard the glide always has somewhere legal to
go.

\subsubsection*{Stage 2 --- the snap (making it balance)}

The glide fixes the endpoints, not the total. In general
$\mathbf{s}\cdot\mathbf{c}\ne d$, so define the shortfall
\[
\Delta = d - \mathbf{s}\cdot\mathbf{c},
\]
and share it out. HardNet would share it \emph{equally}. The anchor shares it
\textbf{in proportion to how much room each channel has} --- because a channel
already pressed against its wall should not be asked to move.

First, which way must channel $i$ go? Its contribution to the total is $s_iP_i$,
so pushing $P_i$ up changes the total by $s_i$. To help close a gap of sign
$\Delta$, channel $i$ must move up when $\Delta s_i > 0$ and down when
$\Delta s_i < 0$. Measure its room \emph{in that direction}:
\[
\text{room}_i =
\begin{cases}
u_i - c_i & \text{if } \Delta s_i > 0 \quad(\text{it must go up}),\\[2pt]
c_i - \ell_i & \text{if } \Delta s_i < 0 \quad(\text{it must go down}).
\end{cases}
\]
Then hand out the shortfall in proportion to that room:
\[
\boxed{\;\mathbf{a} = \mathbf{c}
 + \frac{\Delta}{\sum_m \text{room}_m}\,\bigl(\text{room}_i\, s_i\bigr)_{i=1}^{n}\;}
\]

\paragraph{Check that it balances.} Take the dot product with $\mathbf{s}$ and use
$s_i^2=1$:
\[
\mathbf{s}\cdot\mathbf{a}
= \mathbf{s}\cdot\mathbf{c}
 + \frac{\Delta}{\sum_m \text{room}_m}\sum_i \text{room}_i\,\underbrace{s_i s_i}_{=\,1}
= \mathbf{s}\cdot\mathbf{c} + \Delta\,\frac{\sum_i \text{room}_i}{\sum_m \text{room}_m}
= \mathbf{s}\cdot\mathbf{c} + \Delta = d. \quad\checkmark
\]
The two sums cancel --- that is the only trick, and it works for \emph{any}
positive weights, not just room. Room is the choice we made.

\paragraph{Check that it stays in the box.} Channel $i$ moves by
$|\Delta|\cdot\text{room}_i / \sum_m\text{room}_m$, which is at most
$\text{room}_i$ provided
\[
|\Delta| \;\le\; \sum_m \text{room}_m ,
\]
and $\text{room}_i$ is exactly the distance to the wall it is moving toward. So
as long as the fleet collectively has enough room, no channel overshoots. If it
does not, the formula asks for more than exists, the final clip bites, and balance
is lost --- the same corner case HardNet reports as \texttt{bal\_short}. Neither
head can conjure capacity that is not there.

\subsubsection*{A number, and the contrast with HardNet}

Three generators, $\mathbf{s}=(1,1,-1)$, everything in $[0,10]$, $d=10$. Say
$\mathbf{P}_{\text{prev}}=(4,2,1)$, $\mathbf{p}_R=(10,5,1)$, and $k=2$ moves left.

\emph{Glide:}
\[
\mathbf{c} = (4,2,1) + \frac{(10,5,1)-(4,2,1)}{3} = (4,2,1)+(2,1,0) = (6,3,1),
\qquad \mathbf{s}\cdot\mathbf{c} = 6+3-1 = 8 .
\]

\emph{Snap:} $\Delta = 10-8 = +2$. Generators 1 and 2 have $\Delta s_i>0$ so they
go up; generator 3 has $s_3=-1$ so $\Delta s_3<0$ and it goes down:
\[
\text{room} = \bigl(\underbrace{10-6}_{4},\ \underbrace{10-3}_{7},\ \underbrace{1-0}_{1}\bigr),
\qquad \sum\text{room} = 12,
\qquad \frac{\Delta}{\sum\text{room}} = \frac{2}{12} = \frac16 .
\]
\[
\mathbf{a} = (6,3,1) + \tfrac16\,(4\cdot 1,\ 7\cdot 1,\ 1\cdot(-1))
= \left(6+\tfrac23,\ 3+\tfrac76,\ 1-\tfrac16\right) \approx (6.667,\ 4.167,\ 0.833).
\]
Check: $6.667+4.167-0.833 = 10$ \checkmark

\emph{What HardNet would have done} with the same shortfall: share it equally,
$\lambda = \Delta/\|\mathbf{s}\|^2 = 2/3$, giving $(6.667,\ 3.667,\ 0.333)$ ---
also balanced.

\begin{center}
\begin{tabular}{lccc}
\toprule
 & generator 1 & generator 2 & generator 3 \\
\midrule
room available & $4$ & $7$ & $1$ \\
anchor's move (by room) & $+0.667$ & $+1.167$ & $-0.167$ \\
equal share (HardNet) & $+0.667$ & $+0.667$ & $-0.667$ \\
\bottomrule
\end{tabular}
\end{center}

Generator 2 is roomy, so the anchor leans on it. Generator 3 has only $1$ MW of
room, so the anchor barely touches it while the equal share shoves it $0.667$ of
the way to its floor. That is the entire argument for the weighting: \emph{spend
the shortfall where there is space to spend it.}

\paragraph{The cost of that choice.} The anchor is now a point we picked, not a
point the problem forced on us. Change the weights and the anchor moves; change
the anchor and RAYEN's final answer moves with it. HardNet's ``equal share'' looks
cruder here, but it is not a preference --- it is what minimising
$\|\mathbf{P}-\mathbf{v}\|$ produces, as Section 2 proved. This is the real
difference between the two methods, and it is a genuine trade: RAYEN can encode
knowledge that closeness alone would miss, at the price of that knowledge being an
assumption someone has to defend.

\subsection*{Step 2: aim at the guess}

The direction from anchor to guess is $\mathbf{r} = \mathbf{v}-\mathbf{a}$, and
the points along the way are $\mathbf{a}+t\,\mathbf{r}$ for $t\in[0,1]$.

\subsection*{Step 3: why we cannot walk that way}

Check Rule 1 on the path:
\[
\mathbf{s}\cdot(\mathbf{a}+t\,\mathbf{r})
= \underbrace{\mathbf{s}\cdot\mathbf{a}}_{=\,d} + t\,(\mathbf{s}\cdot\mathbf{r})
= d + t\,(\mathbf{s}\cdot\mathbf{r}).
\]
For this to equal $d$ at \emph{every} $t$ we need
\[
\boxed{\;\mathbf{s}\cdot\mathbf{r} = 0\;}
\]
Walking along $\mathbf{r}$ must not change the total --- it may change
\emph{which} generator supplies the power, never \emph{how much}.

\subsection*{Step 4: fixing the direction --- this is also a projection}

We need the part of $\mathbf{r}$ that is perpendicular to $\mathbf{s}$. That is
exactly the split from Section 2, and discarding the $\mathbf{s}$-component is the
projection of $\mathbf{r}$ onto the plane's parallel directions
$\mathbf{s}^{\perp}$:
\[
\boxed{\;\mathbf{r}' \;=\; \Pi_{\mathbf{s}^{\perp}}(\mathbf{r})
\;=\; \mathbf{r} - \frac{\mathbf{s}\cdot\mathbf{r}}{\|\mathbf{s}\|^{2}}\,\mathbf{s}
\;=\; \mathbf{r} - \frac{\mathbf{s}\cdot\mathbf{r}}{6}\,\mathbf{s}\;}
\]
Check it works:
\[
\mathbf{s}\cdot\mathbf{r}'
= \mathbf{s}\cdot\mathbf{r} - \frac{\mathbf{s}\cdot\mathbf{r}}{6}\,\underbrace{(\mathbf{s}\cdot\mathbf{s})}_{=\,6}
= \mathbf{s}\cdot\mathbf{r} - \mathbf{s}\cdot\mathbf{r} = 0 .\quad\checkmark
\]

\subsection*{Step 5: how far can we walk?}

Coordinate $i$ travels along $a_i + t\,r_i'$. It reaches its ceiling when
\[
a_i + t\,r_i' = u_i \;\Longrightarrow\; t = \frac{u_i-a_i}{r_i'}
\qquad (\text{relevant only if } r_i'>0),
\]
and its floor when
\[
a_i + t\,r_i' = \ell_i \;\Longrightarrow\; t = \frac{a_i-\ell_i}{-r_i'}
\qquad (\text{relevant only if } r_i'<0).
\]
Stop at the first wall any coordinate reaches, and never overshoot the guess:
\[
\boxed{\;t^\star = \min\Bigl(1,\ \min_i\{\text{all such } t\}\Bigr)\;}
\]

\subsection*{Step 6: the answer}
\[
\boxed{\;\mathbf{P} = \mathbf{a} + t^\star\,\mathbf{r}'\;}
\]
Rule 1 holds because $\mathbf{s}\cdot\mathbf{r}'=0$; Rule 2 because $t^\star$
stops before any wall.

\paragraph{What RAYEN is \emph{not}.} It is not $\Pi_C(\mathbf{v})$. It is a point
on a ray that starts at $\mathbf{a}$, and moving $\mathbf{a}$ moves the answer.

\section{Approach B --- HardNet}

HardNet computes $\Pi_C(\mathbf{v})$ exactly:
\[
\mathbf{P} \;=\; \arg\min_{\mathbf{P}\in C}\|\mathbf{P}-\mathbf{v}\|^{2},
\qquad C = C_0\cap B .
\]
Section 2 solved this without the box. Now we add it back.

\subsection*{Step 1: the shape of the answer}

\begin{quote}
\textbf{Claim.} The solution is
$\;\mathbf{P} = \operatorname{clip}(\mathbf{v}+\lambda\mathbf{s},\ \bm{\ell},\ \mathbf{u})\;$
for a \emph{single} number $\lambda$.
\end{quote}

\noindent\textbf{Proof.} Write $\delta_i = P_i - v_i$. The problem is
\[
\min \sum_i \delta_i^{2}
\quad\text{subject to}\quad
\sum_i s_i\delta_i = c,
\qquad
\ell_i - v_i \le \delta_i \le u_i - v_i .
\]
Call a coordinate \textbf{free} if $\delta_i$ is strictly between its bounds, and
\textbf{stuck} otherwise.

Take two free coordinates $i$ and $j$ and nudge them against each other:
\[
\delta_i \to \delta_i + \varepsilon s_i,
\qquad
\delta_j \to \delta_j - \varepsilon s_j .
\]
For small enough $\varepsilon$ of either sign this stays inside the bounds, because
both were strictly inside. It also keeps Rule 1, since
\[
\Delta\Bigl(\sum_k s_k\delta_k\Bigr)
= s_i(\varepsilon s_i) - s_j(\varepsilon s_j)
= \varepsilon\,(s_i^{2}-s_j^{2})
= \varepsilon\,(1-1) = 0 \qquad\text{since } s_i^{2}=s_j^{2}=1 .
\]
The cost changes by
\[
(\delta_i+\varepsilon s_i)^{2} + (\delta_j-\varepsilon s_j)^{2} - \delta_i^{2} - \delta_j^{2}
= 2\varepsilon\,(s_i\delta_i - s_j\delta_j) + \varepsilon^{2}\underbrace{(s_i^{2}+s_j^{2})}_{=\,2}.
\]
At a minimum this must be $\ge 0$ for \emph{both} signs of $\varepsilon$. Divide by
$\varepsilon$ and let $\varepsilon\to0$ from each side: the $\varepsilon^{2}$ term
disappears and we are left with
\[
s_i\delta_i - s_j\delta_j \ge 0
\quad\text{and}\quad
s_i\delta_i - s_j\delta_j \le 0
\qquad\Longrightarrow\qquad
s_i\delta_i = s_j\delta_j .
\]
So \emph{all free coordinates share one common value} of $s_i\delta_i$; call it
$\lambda$. Multiplying by $s_i$ and using $s_i^{2}=1$,
\[
\delta_i = \lambda s_i
\qquad\Longrightarrow\qquad
P_i = v_i + \lambda s_i \quad\text{for every free } i .
\]
A stuck coordinate sits at whichever bound it hit, which is precisely what
$\operatorname{clip}$ returns. Hence
$\mathbf{P}=\operatorname{clip}(\mathbf{v}+\lambda\mathbf{s},\bm{\ell},\mathbf{u})$.
\hfill$\blacksquare$

\paragraph{Sanity check.} With no bounds active every coordinate is free, so
$\mathbf{P}=\mathbf{v}+\lambda\mathbf{s}$ --- Section 2 exactly.

\subsection*{Step 2: finding $\lambda$}

$\lambda$ is whatever makes Rule 1 hold. Define
\[
g(\lambda) = \mathbf{s}\cdot\operatorname{clip}(\mathbf{v}+\lambda\mathbf{s},\ \bm{\ell},\ \mathbf{u})
\]
and solve $g(\lambda)=d$.

\begin{quote}
\textbf{Claim.} $g$ is non-decreasing, so the solution is unique.
\end{quote}

\noindent\textbf{Proof.} Look at one term,
$s_i\operatorname{clip}(v_i+\lambda s_i,\ \ell_i,\ u_i)$, and take $\lambda_1<\lambda_2$.

\emph{If $s_i=+1$:} the inside goes $v_i+\lambda_1 \le v_i+\lambda_2$. Clipping is
non-decreasing (it is the identity in the middle and constant outside), so
$\operatorname{clip}(v_i+\lambda_1)\le\operatorname{clip}(v_i+\lambda_2)$.
Multiplying by $+1$ keeps the inequality.

\emph{If $s_i=-1$:} the inside goes $v_i-\lambda_1 \ge v_i-\lambda_2$, so
$\operatorname{clip}(v_i-\lambda_1)\ge\operatorname{clip}(v_i-\lambda_2)$.
Multiplying by $-1$ reverses it, giving $\le$ again.

Every term is non-decreasing, so $g$ is. A non-decreasing function meets the level
$d$ on at most one interval, so $\lambda$ is essentially unique and can be found by
bisection. \hfill$\blacksquare$

\subsection*{Step 3: the answer}
\[
\boxed{\;\mathbf{P} = \Pi_C(\mathbf{v}) = \operatorname{clip}(\mathbf{v}+\lambda^\star\mathbf{s},\ \bm{\ell},\ \mathbf{u})\;}
\]

\paragraph{Shortcut instead of bisection.} Once the stuck set is known, the free
coordinates satisfy $P_i=v_i+\lambda s_i$, so substituting into Rule 1 and solving
gives $\lambda$ in closed form. With $F$ the free set and $m=|F|$,
\[
d = \sum_{i\notin F} s_iP_i + \sum_{i\in F} s_i(v_i+\lambda s_i)
  = \sum_{i\notin F} s_iP_i + \sum_{i\in F} s_iv_i + \lambda m
\;\Longrightarrow\;
\lambda = \frac{d-\sum_{i\notin F}s_iP_i-\sum_{i\in F}s_iv_i}{m}.
\]
This is the same expansion as Section 2 with the stuck coordinates moved across.

\section{A worked example with three generators}

Take $n=3$, $\mathbf{s}=(1,1,-1)$, $d=10$, every generator in $[0,10]$, and
\[
\mathbf{v}=(12,\;1,\;1).
\]
It fails twice: $v_1=12$ is above its ceiling, and
$\mathbf{s}\cdot\mathbf{v}=12+1-1=12$, overshooting $d=10$ by $2$. Here
$\|\mathbf{s}\|^{2}=3$.

\subsection*{HardNet}

Ignoring the box, $\lambda = \dfrac{10-12}{3} = -\dfrac{2}{3}$ and
\[
\mathbf{v}+\lambda\mathbf{s} = \left(12-\tfrac23,\ 1-\tfrac23,\ 1+\tfrac23\right)
= (11.33,\ 0.33,\ 1.67).
\]
The first coordinate exceeds $10$, so it clips: $P_1=10$, stuck. The free set is
$\{2,3\}$, $m=2$, and the shortcut gives
\[
\lambda = \frac{10 - (1)(10) - \bigl[(1)(1)+(-1)(1)\bigr]}{2} = \frac{0}{2} = 0,
\qquad
\mathbf{P} = (10,\ 1,\ 1).
\]
Check $10+1-1=10$ \checkmark, all inside $[0,10]$ \checkmark, and
$\|\mathbf{P}-\mathbf{v}\| = \|(-2,0,0)\| = 2$.

\subsection*{RAYEN}

Anchor $\mathbf{a}=(5,5,0)$, legal since $5+5-0=10$. Then
\[
\mathbf{r} = (7,\,-4,\,1), \qquad \mathbf{s}\cdot\mathbf{r} = 2,
\qquad
\mathbf{r}' = \mathbf{r} - \tfrac{2}{3}\mathbf{s}
= \left(\tfrac{19}{3},\ -\tfrac{14}{3},\ \tfrac{5}{3}\right),
\qquad \mathbf{s}\cdot\mathbf{r}' = 0\ \checkmark
\]
Walls:
\[
t_1=\frac{10-5}{19/3}=\frac{15}{19}\approx0.789,
\qquad
t_2=\frac{5-0}{14/3}=\frac{15}{14}\approx1.071,
\qquad
t_3=\frac{10-0}{5/3}=6 .
\]
So $t^\star=\tfrac{15}{19}$ and
\[
\mathbf{P} = (5,5,0)+\tfrac{15}{19}\left(\tfrac{19}{3},-\tfrac{14}{3},\tfrac{5}{3}\right)
= \left(10,\ \tfrac{25}{19},\ \tfrac{25}{19}\right) \approx (10,\ 1.316,\ 1.316),
\]
with $\|\mathbf{P}-\mathbf{v}\|\approx 2.049$.

\subsection*{Comparing the two}

Both are legal. HardNet's is closer ($2.000$ against $2.049$) --- necessarily so,
since being closest is the question it was asked. RAYEN's answer would move if we
chose a different anchor; HardNet's would not.

\section{Side by side}

\begin{center}
\begin{tabular}{lll}
\toprule
 & \textbf{RAYEN} & \textbf{HardNet} \\
\midrule
what is projected & the direction $\mathbf{r}$, onto $\mathbf{s}^{\perp}$ & the point $\mathbf{v}$, onto $C$ \\
output equals & $\mathbf{a}+t^\star\Pi_{\mathbf{s}^{\perp}}(\mathbf{r})$ & $\Pi_C(\mathbf{v})$ \\
needs a legal starting point & yes, the anchor & no \\
how the shortfall is shared & in proportion to room & equally, $\lambda$ MW each \\
that sharing rule is & a modelling choice & forced by ``closest'' \\
answer depends on & anchor \emph{and} guess & guess only \\
if the guess is already legal & generally moves it & returns it unchanged \\
extra model outputs & a direction and a step size & the six values \\
finding the scalar & one formula per wall, then a min & bisection on $g(\lambda)=d$ \\
\bottomrule
\end{tabular}
\end{center}

\paragraph{The one-sentence difference.} RAYEN says \emph{begin somewhere legal and
move toward the model as far as the rules allow}; HardNet says \emph{of everything
legal, take whatever is nearest to the model}. Both use projection --- RAYEN on a
direction, HardNet on the point itself --- and only HardNet's answer is free of a
choice we made.

\end{document}
